Biodiversity monitoring using camera traps has become an essential tool for ecological studies. However, manually analyzing the vast volume of data generated by these devices remains a significant challenge. In recent years, deep learning-based methods have been widely used to automate tasks such as filtering empty images, recognizing species, and identifying animal behavior. Despite these advances, such approaches often rely on large volumes of annotated data and face generalization limitations in real-world scenarios. This proposal focuses on investigating the use of multimodal machine learning models—specifically vision-language models—to automate camera trap workflows in low-supervision settings. The first stage of the work investigates the ability of multimodal models to perform three fundamental tasks: filtering empty images, classifying species, and identifying behavior. These tasks are analyzed in zero-shot and few-shot scenarios, also considering inference based on image sequences. Results indicate that multimodal models can achieve competitive performance without task-specific training, although their effectiveness varies depending on the task, model architecture, and use of temporal context. In a second stage, this work proposes and examines a multimodal fusion strategy for species classification in a zero-shot scenario. The approach combines visual similarity obtained via the CLIP model with textual similarity derived from descriptions generated by language models. Results show that integrating visual and textual signals can improve classification accuracy in most evaluated scenarios, although gains depend on the dataset and the species. Finally, this proposal presents WildBind, a multimodal architecture that integrates empty image filtering, species classification, and behavior identification as a single information retrieval problem based on a shared ecological ontology. Wildbind combines an arbitrary number of visual, textual, and contextual evidence sources using a learned, uncertainty-aware fusion mechanism. Overall, this proposal contributes empirical and methodological evidence regarding the use of pre-trained multimodal models in camera-trap workflows, while also presenting a unified architecture as a low-supervision alternative capable of simultaneously addressing the three tasks under investigation.

% {\setlength{\parindent}{0cm} % \textbf{Keywords}: Deep neural networks, camera trap data, automated species % identification, individual animal counting, wildlife monitoring % }